% GENERATED FILE. Edit canonical lessons, then run npm run generate:books.
% canonical-source-hash: fnv1a64:82637cb5b725f38b
% canonical-lessons: ZH-C291-caifang, ZH-C291-faming, ZH-W291-fa-send, ZH-C291-shijian, ZH-C291-tufa, ZH-C291-xuanbu

\chapter{Something New}
\label{ch:zh-something-new}
\hlchaptermodality{291}

\begin{chapteropening}
\textbf{By the end of this chapter:} \emph{I can say someone was interviewed, something was invented, an event broke out and news was announced.}

Complete the last lesson of chapter 291: I can say someone was interviewed, something was invented, an event broke out and news was announced.
\end{chapteropening}

\section[cǎifǎng]{cǎifǎng — to interview someone for the news}

\label{lesson:ZH-C291-caifang}



\begin{quote}
Before the new one: say the Mandarin for the presenter of a show, then the Mandarin for what everyone is talking about.
\end{quote}

\subsection*{You'll want to know: cǎifǎng}
\emph{cǎifǎng} — \textquotedblleft{}to interview someone for the news\textquotedblright{}.

\begin{quote}
\emph{jìzhě cǎifǎng le wǒmen} — A reporter interviewed us.
\end{quote}

Its characters come later; for now it is a word for the ear.

Two third tones meet in this word: cǎi is said with a rising tone before fǎng, and the spelling keeps both marks.

\begin{cousinweb}
\emph{cǎifǎng} is built from \emph{cǎi} 'to pick' and \emph{fǎng} 'to ask about', both inherited from Middle Chinese, \emph{tshojX} and \emph{phjangH} (Guangyun, in Baxter's transcription), which glosses the first 'to take' and the second 'to consult': to gather news by asking. \emph{phjangH} would regularly give fàng; the dipping tone of fǎng is irregular.
\end{cousinweb}

\begin{grammarlens}[title={What you've built}]
One more word for telling what happened.
\end{grammarlens}

\subsection*{Your turn}
\begin{itemize}
\raggedright
  \item \emph{Say it:} \emph{cǎifǎng}
  \item \emph{Say it:} \emph{cǎifǎng}, once more
  \item \emph{Say it:} \emph{jìzhě cǎifǎng le wǒmen} — A reporter interviewed us.
  \item \emph{Recall:} say the Mandarin for the presenter of a show, then the Mandarin for what everyone is talking about, then say \emph{cǎifǎng} again
\end{itemize}

\begin{tcolorbox}[breakable,colback=teal!4,colframe=teal!35!black,title={Before you move on}]
What does cǎifǎng mean? (\textquotedblleft{}To interview someone for the news\textquotedblright{}.) Say it once more.
\end{tcolorbox}
\section[\texorpdfstring{\zh{发明} (fāmíng)}{发明 (fāmíng)}]{\zh{发明} (fāmíng) — to invent; an invention}

\label{lesson:ZH-C291-faming}



\begin{quote}
Before the new one: say the Mandarin for what everyone is talking about, then the Mandarin for to interview someone for the news.
\end{quote}

\subsection*{You'll want to know: \zh{发明}}
\textbf{\zh{发明}} — \emph{fāmíng} — \textquotedblleft{}to invent; an invention\textquotedblright{}.

\begin{quote}
\textbf{\zh{他发明了一个小东西}} \emph{tā fāmíng le yí ge xiǎo dōngxi} — He invented a little gadget.
\end{quote}

The new character, in pieces: \zh{发} — five strokes in three small groups: a short stroke that drops and turns to the right, and a long sweep down to the left from the top; then a fork of two strokes crossing below; then a dot at the top right. The next lesson puts it in your hand.

\begin{cousinweb}
\textbf{\zh{发明}} (\emph{fāmíng}) is built from \zh{发} 'to send out' and \zh{明} 'bright', both inherited from Middle Chinese, \emph{pjot} and \emph{mjaeng} (Guangyun, in Baxter's transcription), which glosses the first 'to start' and 'to make clear' and the second 'light' and 'clear': to bring something to light.
\end{cousinweb}

\begin{grammarlens}[title={What you've built}]
One more word for telling what happened.
\end{grammarlens}

\subsection*{Your turn}
\begin{itemize}
\raggedright
  \item \emph{Say it:} \emph{fāmíng}
  \item \emph{Say it:} \emph{fāmíng}, once more
  \item \emph{Say it:} \emph{tā fāmíng le yí ge xiǎo dōngxi} — He invented a little gadget.
  \item \emph{Recall:} say the Mandarin for what everyone is talking about, then the Mandarin for to interview someone for the news, then say \emph{fāmíng} again
\end{itemize}

\begin{tcolorbox}[breakable,colback=teal!4,colframe=teal!35!black,title={Before you move on}]
What does \zh{发明} mean? (\textquotedblleft{}To invent; an invention\textquotedblright{}.) Say it once more.
\end{tcolorbox}
\section[\texorpdfstring{\zh{发} (fā)}{发 (fā)}]{\zh{发} — a bent stroke, a sweep, a fork and a dot}

\label{lesson:ZH-W291-fa-send}



\begin{quote}
The word first, then the character inside it.

\begin{itemize}
\raggedright
  \item \emph{Recall:} say \emph{fāmíng} — \textbf{R1}, one lesson back, before it decays
\end{itemize}

\begin{quote}
\zh{突} \zh{明}
\end{quote}

\emph{Read it:} each one aloud, with nothing said first — \emph{tū}, \emph{míng}
\end{quote}

\subsection*{Script — a bent stroke, a sweep, a fork and a dot}
\hlblockfigure{\detokenize{figures/ZH-W291-fa-send-filmstrip.pdf}}{How \zh{发} is written, stroke by stroke}

\begin{quote}
\zh{发}
\end{quote}

\textbf{One shape, three small groups.} Count the groups, not the strokes.

\textbf{The frame, two strokes.} A short stroke that drops and turns to run across to the right, rising a little. Then a long sweep from the top middle, down to the lower left, crossing it.

\textbf{The fork, two strokes.} A bar that turns into a sweep down to the lower left. Then a long stroke falling to the right from under the bar, out to the foot.

\textbf{The dot, one stroke}, high on the right.

2 + 2 + 1 = five.

\subsection*{Your turn}
\begin{itemize}
\raggedright
  \item \emph{Write it:} \zh{发}, all five, with the model in front of you
  \item \emph{Cover:} write it again from memory, and count the pieces, not the strokes
\end{itemize}

\begin{tcolorbox}[breakable,colback=teal!4,colframe=teal!35!black,title={Before you move on}]
How many strokes in \zh{发}? (\textbf{Five}.) Which stroke comes first? (\textbf{The short one that drops and turns right}.) What comes last? (\textbf{The dot, high on the right}.)

Source: \href{https://raw.githubusercontent.com/chanind/hanzi-writer-data/68d10a4b21150cae5e1ebbd223eed289cf32d90c/data/\%E5\%8F\%91.json}{Hanzi Writer Data \zh{发}, pinned stroke paths and medians}.
\end{tcolorbox}
\section[shìjiàn]{shìjiàn — incident; event (in the news)}

\label{lesson:ZH-C291-shijian}



\begin{quote}
Before the new one: say the Mandarin for to interview someone for the news, then the Mandarin for to invent something.
\end{quote}

\subsection*{You'll want to know: shìjiàn}
\emph{shìjiàn} — \textquotedblleft{}incident; event (in the news)\textquotedblright{}.

\begin{quote}
\emph{zhège shìjiàn fāshēng zài Běijīng} — The incident happened in Beijing.
\end{quote}

Its characters come later; for now it is a word for the ear.

\begin{cousinweb}
\emph{shìjiàn} is built from \zh{事} 'a matter' and \emph{jiàn} 'an item', both inherited from Middle Chinese, \emph{dzriH} and \emph{gjenX} (Guangyun, in Baxter's transcription), which glosses the second 'to divide into parts'.
\end{cousinweb}

\begin{grammarlens}[title={What you've built}]
One more word for telling what happened.
\end{grammarlens}

\subsection*{Your turn}
\begin{itemize}
\raggedright
  \item \emph{Say it:} \emph{shìjiàn}
  \item \emph{Say it:} \emph{shìjiàn}, once more
  \item \emph{Say it:} \emph{zhège shìjiàn fāshēng zài Běijīng} — The incident happened in Beijing.
  \item \emph{Recall:} say the Mandarin for to interview someone for the news, then the Mandarin for to invent something, then say \emph{shìjiàn} again
\end{itemize}

\begin{tcolorbox}[breakable,colback=teal!4,colframe=teal!35!black,title={Before you move on}]
What does shìjiàn mean? (\textquotedblleft{}Incident; event (in the news)\textquotedblright{}.) Say it once more.
\end{tcolorbox}
\section[\texorpdfstring{\zh{突发} (tūfā)}{突发 (tūfā)}]{\zh{突发} (tūfā) — to break out suddenly; sudden}

\label{lesson:ZH-C291-tufa}



\begin{quote}
Before the new one: say the Mandarin for to invent something, then the Mandarin for an incident in the news.
\end{quote}

\subsection*{You'll want to know: \zh{突发}}
\textbf{\zh{突发}} — \emph{tūfā} — \textquotedblleft{}to break out suddenly; sudden\textquotedblright{}.

\begin{quote}
\emph{zhè shì yí ge tūfā shìjiàn} — This is a sudden, unexpected event.
\end{quote}

\begin{cousinweb}
\textbf{\zh{突发}} (\emph{tūfā}) is built from \zh{突} 'to burst out' and \zh{发} 'to start', both inherited from Middle Chinese, \emph{thwot} and \emph{pjot} (Guangyun, in Baxter's transcription), which glosses the first 'the look of something coming out'.
\end{cousinweb}

\begin{grammarlens}[title={What you've built}]
One more word for telling what happened.
\end{grammarlens}

\subsection*{Your turn}
\begin{itemize}
\raggedright
  \item \emph{Say it:} \emph{tūfā}
  \item \emph{Say it:} \emph{tūfā}, once more
  \item \emph{Say it:} \emph{zhè shì yí ge tūfā shìjiàn} — This is a sudden, unexpected event.
  \item \emph{Recall:} say the Mandarin for to invent something, then the Mandarin for an incident in the news, then say \emph{tūfā} again
\end{itemize}

\begin{tcolorbox}[breakable,colback=teal!4,colframe=teal!35!black,title={Before you move on}]
What does \zh{突发} mean? (\textquotedblleft{}To break out suddenly; sudden\textquotedblright{}.) Say it once more.
\end{tcolorbox}
\section[xuānbù]{xuānbù — to announce officially}

\label{lesson:ZH-C291-xuanbu}



\begin{quote}
Before the new one: say the Mandarin for an incident in the news, then the Mandarin for to break out suddenly.
\end{quote}

\subsection*{You'll want to know: xuānbù}
\emph{xuānbù} — \textquotedblleft{}to announce officially\textquotedblright{}.

\begin{quote}
\emph{xuéxiào xuānbù míngtiān fàngjià} — The school announced a day off tomorrow.
\end{quote}

Its characters come later; for now it is a word for the ear.

\begin{cousinweb}
\emph{xuānbù} is built from \emph{xuān} 'to make known' and \emph{bù} 'to spread', both inherited from Middle Chinese, \emph{sjwen} and \emph{puH} (Guangyun, in Baxter's transcription), which glosses the first 'to spread' and 'to make clear' and the second 'to spread everywhere'.
\end{cousinweb}

\begin{grammarlens}[title={What you've built}]
One more word for telling what happened.
\end{grammarlens}

\subsection*{Your turn}
\begin{itemize}
\raggedright
  \item \emph{Say it:} \emph{xuānbù}
  \item \emph{Say it:} \emph{xuānbù}, once more
  \item \emph{Say it:} \emph{xuéxiào xuānbù míngtiān fàngjià} — The school announced a day off tomorrow.
  \item \emph{Recall:} say the Mandarin for an incident in the news, then the Mandarin for to break out suddenly, then say \emph{xuānbù} again
\end{itemize}

\begin{tcolorbox}[breakable,colback=teal!4,colframe=teal!35!black,title={Before you move on}]
What does xuānbù mean? (\textquotedblleft{}To announce officially\textquotedblright{}.) Say it once more.
\end{tcolorbox}
